
Our work intersects three research areas: selective prediction for machine learning models, uncertainty quantification in large language models, and experimental methodology in machine learning. We review each area to position our contributions and highlight the gap in experimental design guidelines that our work addresses.

\subsubsection{Selective Prediction and Uncertainty Estimation}

Selective prediction, introduced by Chow (1957) and formalized by El-Yaniv \& Wiener (2010), allows models to abstain from predictions when uncertainty exceeds a threshold. The core challenge is estimating uncertainty without access to ground-truth labels at test time. Hendrycks \& Gimpel (2017) demonstrated that maximum softmax probability provides a surprisingly effective baseline for out-of-distribution detection in image classification, establishing max-probability thresholding as the standard single-forward-pass method. Liang et al. (2018) improved upon this by applying temperature scaling and input perturbations, showing that simple post-processing can enhance uncertainty signals.

However, max-probability captures only the mode of the probability distribution---the single most likely outcome. For autoregressive language models generating token-by-token predictions, this mode-only signal may miss uncertainty patterns in the full distribution. If a model assigns 40\% probability to one token, 30\% to another, and 30\% to a third, max-probability reports only 40\%, discarding information about the distribution's flatness. This motivates alternative uncertainty measures that capture distributional shape.

\subsubsection{Entropy-Based Uncertainty Quantification}

Shannon entropy \citep{Shannon1948}, defined as H = -$\Sigma$ p(x) log p(x), quantifies the unpredictability of a distribution. In machine learning, entropy has been applied to measure model confidence: peaked distributions (high confidence) have low entropy, while flat distributions (high uncertainty) have high entropy. Gal \& Ghahramani (2016) used predictive entropy in Bayesian neural networks with MC Dropout to estimate epistemic uncertainty. Malinin \& Gales (2018) distinguished between distributional uncertainty (entropy of output distribution) and distributional ambiguity (entropy over parameter posteriors), arguing that both are necessary for reliable uncertainty estimates.

For large language models, entropy offers a theoretically appealing alternative to max-probability because it captures multi-modal distributions. If a model is uncertain between multiple plausible continuations, entropy increases even when the top probability remains moderately high. Kuhn et al. (2023) explored semantic entropy for long-form generation, clustering semantically equivalent continuations before computing entropy to better capture meaningful uncertainty. However, these methods have primarily been evaluated on generative tasks, not selective prediction for factual question answering.

The critical gap in this literature is the assumption that entropy can be straightforwardly evaluated against any model on any dataset. Existing work compares entropy-based methods to max-probability baselines, but does not establish experimental prerequisites---specifically, whether the chosen model produces sufficient variance in correctness for correlation-based statistical tests to be valid. Our work fills this gap by identifying model capacity as a binary gate for test validity.

\subsubsection{Model Scaling and Factual Knowledge}

The relationship between model size and factual knowledge capacity is well-established in the scaling laws literature. Kaplan et al. (2020) showed that language model performance follows predictable power laws with respect to model parameters, dataset size, and compute. Brown et al. (2020) demonstrated that GPT-3 (175B parameters) exhibits emergent few-shot learning abilities absent in smaller models, attributing this to the compression of factual knowledge acquired during pre-training. Rae et al. (2021) and Hoffmann et al. (2022) further refined these scaling laws for compute-optimal training, showing that model capacity directly determines the amount of factual knowledge that can be stored.

For factual question answering specifically, Petroni et al. (2019) found that BERT-base models store relational knowledge but perform poorly on complex factual queries. Roberts et al. (2020) showed that closed-book question answering (no retrieval) requires models of at least 10B parameters to achieve competitive performance on TriviaQA. GPT-2 (117M parameters), designed for general language modeling rather than knowledge recall, performs at near-zero accuracy on TriviaQA. This explains our zero-variance correctness outcome, but prior work has not connected model capacity limitations to the validity of uncertainty estimation experiments.

\subsubsection{Experimental Design in Machine Learning}

The machine learning community has increasingly recognized the importance of rigorous experimental methodology. Dwork et al. (2015) formalized fairness definitions and evaluation protocols, establishing that valid fairness tests require specific distributional conditions. Henderson et al. (2018) showed that deep reinforcement learning experiments suffer from high variance due to random seeds and hyperparameter choices, recommending stricter reporting standards. Bouthillier et al. (2021) demonstrated that many published results cannot be reproduced due to insufficient experimental detail.

However, these methodological critiques focus on reproducibility and statistical power, not on prerequisites for test validity. They assume that chosen experimental conditions support the intended statistical analyses---that correlation tests can be computed, that ANOVA assumptions are met, that sample sizes are adequate. Our work reveals a failure mode that invalidates tests entirely: zero-variance outcomes that make correlation mathematically undefined. This is distinct from low power (failing to detect a true effect) or p-hacking (selective reporting). It is a condition where the test itself cannot run, yet current experimental practices provide no guardrails to detect this failure mode at the design stage.

\subsubsection{Positioning Our Work}

Our infrastructure validation builds on standard practices: token probability extraction from language models \citep{Hendrycks2017}, Shannon entropy computation \citep{Shannon1948}, and statistical correlation tests \citep{Spearman1904}. Our contribution is not a novel uncertainty estimation method but rather a methodological insight: selective prediction experiments have dependencies on model capacity that gate the validity of correlation-based statistical tests. We demonstrate this through a concrete failure case---GPT-2 on TriviaQA producing zero-variance correctness---that existing literature would characterize as ''model performs poorly'' rather than ''experiment design is invalid.''

This positions our work as a methodological contribution to experimental design in machine learning uncertainty quantification. We provide empirical validation that infrastructure (entropy extraction, quadrant analysis) works independent of model performance, while identifying model capacity thresholds as a prerequisite for hypothesis testing. Future selective prediction studies must report model capacity explicitly and verify non-zero variance in evaluation metrics before claiming negative results. Our framework for distinguishing infrastructure validation from hypothesis validation applies beyond entropy-based methods to any correlation-based uncertainty evaluation.
